\documentclass[12pt]{article}
\usepackage[utf8]{inputenc}
\usepackage[T1]{fontenc}
\usepackage{hyperref}
\title{Pahanan Agta Language Resources for K-12 Curriculum Development}
\author{Compiled by Claude Sonnet 5.5 (claude-sonnet-5-5) for Andr\'as Kornai}
\date{}
\begin{document}
\maketitle

\section{Introduction}
Pahanan Agta (ISO 639-3 apf; listed as Palanan Agta by Scripture Earth) is an Austronesian language of northeastern Luzon, Philippines, near Palanan. It is closely related to Casiguran Dumagat Agta, Dupaningan Agta and Dinapigue Agta. Search results on Pahanan itself were thin.

\section{Primary Educational Resources}
None found in the search results.

\section{Linguistic and Reference Materials}
\begin{itemize}
\item Survey article on the northeastern Luzon languages, including Pahanan, arguing they form a primary branch of Northern Luzon: \url{https://doi.org/10.1353/ol.2013.0007}
\item Casiguran Dumagat Agta, a closely related language, with documentation by SIL linguists: \url{https://en.wikipedia.org/wiki/Casiguran_Dumagat_Agta}
\end{itemize}
No Pahanan grammar or dictionary was found.

\section{Religious and Community Resources}
Scripture Earth lists scripture recordings (Gospel of Mark and Luke on Bible.is, and Global Recordings audio) for the language: \url{https://scriptureearth.org/00eng.php?iso=apf}

\section{Digital and Audio Resources}
The audio items are those on the Scripture Earth page above.

\section{Curriculum Development Context}
Very little written material exists. A curriculum would need new materials built with the community, possibly drawing on work for related Agta languages.
\end{document}